\newpage
\section{Conclusion}

We have presented a multi‑prover framework for formally verifying the
atomicity of foundational observables in the \Apeiron{} Framework.  By
integrating SymPy, Z3, Lean~4, and Coq, the system can generate,
cross‑verify, and export machine‑checkable proofs of atomicity across
multiple mathematical frameworks.

A rigorous evaluation on a 47‑observable benchmark reveals a clear trust
hierarchy: \textbf{Z3} is the only fully automatic prover that is sound for
Boolean atomicity on both integers and sets; \textbf{Lean~4} and \textbf{Coq}
provide genuine computational proofs for integers, but their backends for
other types are currently placeholders; \textbf{SymPy} serves as a fast
consistency filter but is not sound for atomicity.  We have defined a formal
trust function $\tau$ that assigns each prover a weight of $1$, $\frac{1}{2}$,
or $0$, and we generate certificates that faithfully reflect this hierarchy.

This work demonstrates that pragmatic, verifiable multi‑prover architectures
are feasible today, while openly acknowledging the limitations that remain.
The JSON certificates, together with the trust hierarchy, enable third‑party
auditors to verify the system’s internal claims without trusting the AI
itself—a crucial property for regulatory acceptance.  Future work will focus
on extending the Lean and Coq backends to non‑integer observables,
formalising the translation layer via an abstract syntax tree, and developing
compositional proof rules that scale beyond individual observables.

The open‑source release of \texttt{self\_proving.py} and its accompanying
certificates invites the community to critique, extend, and build upon this
foundation in the pursuit of truly trustworthy artificial intelligence.